\documentclass[10pt,twocolumn]{article}
\usepackage[T1]{fontenc}
\usepackage[utf8]{inputenc}
\usepackage{lmodern}
\usepackage[margin=1.8cm,top=2.4cm,bottom=2cm,columnsep=0.6cm]{geometry}
\usepackage{fancyhdr}
\usepackage{titlesec}
\usepackage{booktabs}
\usepackage{amsmath,amssymb,amsfonts}
\DeclareMathOperator*{\argmin}{arg\,min}
\usepackage{microtype}
\usepackage{xcolor}
\usepackage{mdframed}
\usepackage{parskip}
\usepackage{enumitem}
\usepackage{hyperref}

\definecolor{rulered}{RGB}{180,30,30}
\definecolor{boxgray}{RGB}{245,245,245}
\definecolor{keywordgray}{RGB}{100,100,100}

\pagestyle{fancy}
\fancyhf{}
\fancyhead[L]{\textsc{\small Claude Mag}}
\fancyhead[C]{\small\textit{Machine Learning Research Digest}}
\fancyhead[R]{\small 2026-07-27}
\fancyfoot[C]{\small\thepage}
\renewcommand{\headrulewidth}{0.8pt}

\titleformat{\section}{\normalsize\bfseries\color{rulered}}{}{0pt}{}%
  [\vspace{-4pt}\color{rulered}\rule{\columnwidth}{0.4pt}]
\titlespacing{\section}{0pt}{8pt}{4pt}

\hypersetup{colorlinks=true,urlcolor=rulered,linkcolor=black}

\begin{document}
\setlength{\parindent}{0pt}
\setlength{\parskip}{4pt}

{\small\color{keywordgray}\textbf{Keywords:}
  token compression \textbullet{}
  context distillation \textbullet{}
  LLM efficiency \textbullet{}
  attention mechanisms}\par\vspace{4pt}

{\LARGE\bfseries\raggedright
  Progressive Cramming: Reliable Token Compression
  and What It Reveals}\par\vspace{4pt}

{\large\itshape\raggedright
  A Token-by-Token Compression Protocol Uncovers a Fundamental
  Capacity Limit in LLM Embedding Space---and Traces Downstream
  Task Degradation to Destructive Early-Layer Attention Interactions}\par\vspace{6pt}

{\small\textbf{By} Dmitrii Tarasov, Timofei Lashukov, Elizaveta Goncharova,
  Andrey Kuznetsov}\par\vspace{2pt}

{\small\color{keywordgray}arXiv:2607.21231 \textbullet{} July 23, 2026}
\par\vspace{2pt}\noindent{\color{rulered}\rule{\columnwidth}{0.6pt}}\vspace{6pt}

Token cramming---compressing long input prefixes into a small set of
learned continuous embeddings---promises to let LLMs reference long
contexts at the cost of a handful of extra tokens. Prior work reports
near-perfect reconstruction (99\%+), but fixes both the token budget
and the accuracy threshold, leaving a crucial ambiguity: is residual
error an optimization shortfall, or a fundamental limit of the embedding
space? This paper answers with \textbf{progressive cramming}, a
protocol that grows the compression target token-by-token until failure,
revealing hard content-dependent limits that optimization cannot breach.

\begin{mdframed}[backgroundcolor=boxgray,linecolor=rulered,linewidth=1pt,
    innerleftmargin=6pt,innerrightmargin=6pt,
    innertopmargin=6pt,innerbottommargin=6pt]
\textbf{\small AT A GLANCE}\par\vspace{4pt}
\begin{description}[leftmargin=0pt,labelsep=4pt,itemsep=2pt,parsep=0pt,
    font=\small\bfseries]
  \item[Problem:] Existing token cramming evaluates at fixed budgets
    and thresholds, conflating optimization failure with fundamental
    compression limits.
  \item[Why it matters:] Systems relying on token cramming for
    long-context compression may degrade silently at inference;
    knowing true compression horizons is essential for safe deployment.
  \item[Method:] Progressive cramming grows target prefix
    token-by-token, stopping at reconstruction failure; causal
    attention-knockout interventions locate the degradation source.
  \item[Result:] Crammed embeddings cause 4--9\,pp accuracy drops
    on multiple-choice tasks and near-total collapse on generative
    tasks, traced to early-layer attention interactions.
  \item[Evidence:] Compression horizons range from 28--74\% of
    prefix length depending on content, revealing fundamental limits
    independent of optimizer tuning.
\end{description}
\end{mdframed}

\section{The Big Picture}

The canonical approach to long-context LLM inference is to extend
the context window---but this incurs quadratic attention cost.
Token cramming offers an alternative: compress a prefix of $N$ tokens
into $M \ll N$ learned embedding vectors, insert them into the
context, and let the model proceed as if the prefix were present.
Studies have shown these crammed embeddings can achieve $>99\%$ token-level
reconstruction accuracy, fueling optimism about practical deployment.

This paper asks a harder question: does reconstruction accuracy predict
task performance? And is there a hard ceiling on how many tokens can
be crammed, independent of how long you optimize? The answers are
``no'' and ``yes'', respectively, with implications for any deployment
of cramming-based compression.

\section{Why Existing Approaches Fall Short}

Standard cramming evaluates quality at a fixed compression ratio (e.g.,
compress 512 tokens into 32 embeddings) and declares success when
reconstruction accuracy exceeds 99\%. This methodology has two gaps:

\textbf{Conflation of failure modes.} If reconstruction fails at
a given ratio, it is unclear whether (a) more optimization steps would
succeed or (b) the information exceeds the embedding's representational
capacity. Distinguishing these requires a protocol that rules out
optimization failure as an explanation.

\textbf{Reconstruction $\neq$ task performance.} A crammed embedding
that faithfully reconstructs a prefix token-by-token may still disrupt
downstream task performance if the embedding interferes with the model's
internal processing---even when the original prefix is also provided
in context.

\section{Core Mechanism}

Progressive cramming resolves the first gap by progressively growing
the target prefix from length 1 to length $N$, applying a fixed
optimization budget at each step. At length $n$, it finds the best
embedding for the first $n$ tokens. At some critical length $n^*$,
optimization fails to achieve a fixed reconstruction threshold within
the budget. Because the budget is held constant throughout, failure
at $n^*$ cannot be attributed to insufficient compute---it reflects
a genuine representational limit.

A key structural finding is that successful progressive trajectories
$\{e^*_1, e^*_2, \ldots, e^*_{n^*}\}$ lie in a low-dimensional manifold
in embedding space. PCA on these trajectories shows 80\% of variance
explained by 5 principal components---far fewer than the ambient
embedding dimension of 4096---suggesting that token cramming exploits
a structured low-dimensional code in the embedding space.

\section{Technical Formulation}

Let the LLM $f_\theta$ be frozen with token vocabulary $\mathcal{V}$.
Given prefix $x = (x_1, \ldots, x_N) \in \mathcal{V}^N$, the
cramming objective finds an embedding $e \in \mathbb{R}^d$:

\[
e^*_n = \argmin_{e \in \mathbb{R}^d}
  \mathcal{L}_\text{recon}(e, x_{1:n}), \quad
\mathcal{L}_\text{recon} = \sum_{i=1}^n \ell_\mathrm{CE}(f_\theta(e, x_{<i}), x_i)
\]

subject to an optimization budget of $B$ gradient steps. Progressive
cramming solves this sequence for $n = 1, 2, \ldots, N$, stopping at:

\[
n^* = \max\{n : \mathcal{L}_\text{recon}(e^*_n, x_{1:n}) \leq \tau\}
\]

The \textbf{compression horizon} is $\rho = n^*/N \in (0, 1]$: the
fraction of the prefix representable within budget $B$.

\section{Learning or Inference Procedure}

After establishing compression horizons, the paper tests whether
crammed embeddings $e^*_{n^*}$ degrade task performance by prepending
them to contexts where the full original prefix is also available:

\begin{enumerate}[itemsep=2pt,parsep=0pt]
  \item Insert $e^*$ as a soft prompt token before the original prefix.
  \item Query the model on multiple-choice and open-ended generation tasks.
  \item Compare accuracy to a baseline where no crammed embedding
    is inserted (original prefix only).
\end{enumerate}

To locate the source of degradation, causal attention-knockout
interventions systematically ablate attention edges between $e^*$
and each layer's query/key/value projections, measuring which ablation
most recovers task performance.

\section{What the Guarantee Says}

There is no guarantee that a crammed embedding at the compression horizon
$n^*$ will not degrade task performance. The results show the opposite:
even maximally compressed embeddings (at the representational limit)
cause consistent accuracy drops. The horizon tells practitioners how
much of a prefix the embedding \emph{encodes}---but not how much it
\emph{transmits usefully to the model at inference time}.

\section{Experimental Findings}

Experiments use GPT-2-XL (1.5B), LLaMA-2-7B, and Mistral-7B as
backbone models, with prefixes from C4 (general text), PubMed (medical),
and HumanEval (code).

\begin{center}
\small
\begin{tabular}{lcc}
\toprule
Task Type & No Crammed Embed. & + Crammed Embed. \\
\midrule
MMLU (4-choice) & 63.4 & 56.7 \\
HellaSwag (4-choice) & 74.1 & 67.8 \\
OpenGen (correctness) & 51.3 & 18.2 \\
Code completion & 42.7 & 11.4 \\
\bottomrule
\end{tabular}
\end{center}

Compression horizons (fraction of 512-token prefixes representable):
General text: 68\%, Medical text: 55\%, Code: 28\%.
Code is hardest to compress, consistent with the view that
low-entropy, syntactically precise sequences exhaust embedding capacity fastest.

\section{Ablations and Interpretation}

\textbf{Attention knockout:} Ablating attention edges from $e^*$ in
layers 1--4 recovers 61\% of the multiple-choice accuracy drop.
Ablating in layers 12--16 recovers only 8\%.
Early layers are the primary site of destructive interaction.

\textbf{Prefix vs.\ no prefix:} Inserting $e^*$ with the original prefix
still degrades performance. Inserting $e^*$ without the original prefix
is catastrophic. This rules out the hypothesis that the embedding
``blocks'' the prefix from the model; rather, the embedding actively
interferes with early attention distributions.

\textbf{Budget sensitivity:} Doubling $B$ from $10^3$ to $2 \times 10^3$
steps increases the compression horizon by only 3--6\%, confirming that
the horizon is not a function of insufficient optimization.

\textbf{Dimension sensitivity:} Increasing embedding dimension from
$d = 4096$ to $d = 8192$ raises horizons by 11\% on average,
suggesting that embedding dimensionality is a partial bottleneck,
though the gains plateau rapidly.

\section*{Reference}

Dmitrii Tarasov, Timofei Lashukov, Elizaveta Goncharova, Andrey Kuznetsov.
\textbf{Progressive Cramming: Reliable Token Compression and What It Reveals}.
arXiv:2607.21231, July 23, 2026.
\url{https://arxiv.org/abs/2607.21231}

\end{document}
